\section{Arquitetura em Camadas}

\hspace{0.5cm} A arquitetura em camadas da aplicação segue princípios consolidados de engenharia de software,
integrando padrões de design amplamente adotados na indústria. A implementação fundamenta-se nos seguintes padrões:

\begin{enumerate}
  \item \textbf{Padrões de Design Clássicos}: Os padrões \textit{Template Method}, \textit{Strategy},
  \textit{Decorator} e \textit{Dependency Injection} são classificados como padrões de design estruturais
  e comportamentais conforme definido por \cite{gamma2009design}.

  \item \textbf{Padrões Empresariais}: O padrão \textit{Repository} e \textit{Data Transfer Object (DTO)}
  são padrões de aplicação empresarial descritos por \cite{fowler2015patterns} que facilitam a separação
  entre lógica de negócio e camadas de persistência.

  \item \textbf{Service Layer}: O padrão de organização \textit{Service Layer} concentra a lógica de negócio
  em uma camada dedicada, seguindo arquiteturas em camadas conforme descrito por \cite{hohpe2015enterprise}
  e \cite{evans2014domain}.
\end{enumerate}

\subsection{Domínio}

\hspace{0.5cm} A camada de domínio encapsula a modelagem das entidades de negócio e suas regras de validação.
Esta implementação separa explicitamente as responsabilidades através do padrão DTO, Repository e Service Layer,
garantindo que validações sejam aplicadas em múltiplas camadas e que a lógica de persistência não se misture
com a lógica de negócio.

\subsubsection{Estrutura Padronizada}

Cada domínio é implementado em um arquivo único que consolida três representações complementares da mesma entidade:

\begin{itemize}
  \item \textbf{Schema}: Classe Pydantic que valida dados de entrada para operações de criação.
  Define tipos, restrições de tamanho (\texttt{max\_digits}, \texttt{decimal\_places}), valores padrão e enumerações.
  Não inclui campos gerenciados pelo sistema como \texttt{id} ou \texttt{created\_at}.
  É um DTO que transfere dados entre a API HTTP e as camadas internas.

  \item \textbf{UpdateSchema}: Classe Pydantic com todos os campos como \texttt{Optional},
  permitindo atualizações parciais sem requerer a entidade completa. Esta separação previne validações
  que rejeitariam requisições legítimas com campos ausentes. Também funciona como DTO para operações de modificação.

  \item \textbf{Model}: Classe SQLModel que utiliza o decorador \texttt{table=True}, mapeando a entidade diretamente para tabelas SQL.
  Inclui campos de auditoria automática (\texttt{created\_at}, \texttt{updated\_at}, \texttt{deleted\_at}) para implementar soft delete,
  UUIDs v7 como chaves primárias e relacionamentos através de \texttt{Relationship} com suporte a cardinalidade muitos-para-muitos
  via modelos de ligação. Não contém lógica de persistência, apenas representação de dados.
\end{itemize}

\subsubsection{Enumerações}

Para domínios com estados ou tipologias bem-definidas, enumerações Pydantic/SQLModel são utilizadas.
Isso assegura type-safety em tempo de compilação (estática) e validação em tempo de execução, além de documentar
os valores permitidos de forma explícita. Exemplos incluem:

\begin{itemize}
  \item \texttt{RoleEnum}: Define os papéis de acesso de funcionários (admin, manager, chef, waiter, delivery).
  \item \texttt{StatusEnum}: Estados de um pedido (pending, preparing, in\_route, completed, failed, canceled).
  \item \texttt{ChannelEnum}: Canais de origem de pedidos (app, totem, counter, web).
  \item \texttt{OrderTypeEnum}: Modos de entrega (dine\_in, takeaway, delivery).
  \item \texttt{PaymentMethodEnum}: Métodos de pagamento suportados (cash, credit\_card, pix, vale\_alimentacao).
  \item \texttt{AuditActionEnum}: Ações registradas em auditoria (CREATE, READ, UPDATE, DELETE, LOGIN, LOGOUT, PASSWORD\_CHANGE).
\end{itemize}

\subsubsection{Padrão Repository}

A camada Repository encapsula todas as operações de banco de dados em uma classe genérica
\texttt{BaseRepository<T>}, que implementa CRUD padrão (create, read, update, delete) através de métodos type-safe.
Esta separação mantém consultas SQL e gerenciamento de sessão afastados da lógica de negócio.

Cada domínio pode estender \texttt{BaseRepository} com queries específicas quando necessário. Por exemplo,
\texttt{EmployeeRepository} adiciona \texttt{get\_by\_email()}, permitindo buscar funcionários por email sem duplicar
código de base de dados. O Repository é responsável por:

\begin{itemize}
  \item Gerenciar a sessão SQLModel/SQLAlchemy.
  \item Executar queries de forma segura.
  \item Controlar transações (commit/rollback).
  \item Retornar entidades (Models) ou \texttt{None} em caso de não-encontrado.
\end{itemize}

\subsubsection{Padrão Service}

A camada Service orquestra a transformação entre DTOs (entrada da API) e Models (persistência),
aplicando lógica de negócio, validações e transformações adicionais. Cada domínio possui um Service que:

\begin{itemize}
  \item Consome um Repository para operações de banco de dados.
  \item Valida regras de negócio que vão além da validação estrutural de Schemas.
  \item Aplica transformações de dados (ex: hashing de senhas, busca de endereço por CEP).
  \item Orquestra chamadas a serviços externos ou outros domínios.
  \item Mapeia HTTPExceptions para erros adequados (404, 400, 403).
\end{itemize}

A classe genérica \texttt{BaseService<ModelT, CreateSchemaT, UpdateSchemaT>} implementa operações padrão,
permitindo que Services específicos adicionem métodos customizados. Por exemplo, \texttt{CostumerService.create\_costumer()}
não apenas cria um registro, mas também busca o endereço completo via CEP e aplica hash à senha antes de persistir.

\subsubsection{Fluxo de Dados}

Uma requisição segue o fluxo unidirecional:

\begin{enumerate}
  \item \textbf{Router (APIRouter)}: Recebe a requisição HTTP e valida contra um Schema (DTO).
  \item \textbf{Service}: Recebe o Schema validado, aplica lógica de negócio, cria/atualiza um Model.
  \item \textbf{Repository}: Recebe o Model e persiste no banco de dados via SQLModel.
  \item \textbf{Resposta}: O Model é serializado novamente como JSON e retornado ao cliente.
\end{enumerate}

\subsubsection{Domínios Implementados}

A aplicação organiza suas entidades em sete domínios principais:

\begin{description}
  \item[Audit] Registra todas as operações críticas do sistema (criação, leitura, atualização, exclusão, autenticação).
  Armazena metadados como IP de origem, user agent, status HTTP e identificador do executor, permitindo rastreabilidade completa.

  \item[Auth] Gerencia autenticação e tokens JWT para acesso à API. Separado em domínio próprio para manter segurança em primeiro plano.

  \item[Catalogue] Representa o cardápio de pratos oferecidos. Estabelece relacionamento muitos-para-muitos com \texttt{Inventory}
  através do modelo de ligação \texttt{DishIngredient}, que mapeia quais ingredientes compõem cada prato.

  \item[Costumer] Modela clientes/usuários finais com autenticação, dados de contato, endereço como JSON estruturado e
  sistema de pontos (loyalty). Mantém histórico de pedidos também em JSON para flexibilidade.

  \item[Employee] Define funcionários com papéis de autorização (\texttt{RoleEnum}), dados de endereço e autenticação separada.
  Estrutura preparada para hierarquia de acesso baseada em roles.

  \item[Inventory] Controla o estoque de ingredientes com unidades de medida específicas (\texttt{UnitEnum}: gramas, mililitros, unidades).
  Define quantidades mínimas para alertas de reposição. Conectado ao \texttt{Catalogue} através do modelo de ligação.

  \item[Order] Encapsula o conceito central de negócio: pedidos. Agrupa múltiplas dimensões (origem, tipo, pagamento, status)
  e armazena detalhes de itens como JSON estruturado para flexibilidade, pois cada item pode ter variações dinâmicas.
\end{description}

\subsubsection{Padrões de Persistência}

Todas as entidades implementam soft delete através do campo \texttt{deleted\_at}. Isso preserva integridade referencial
e permite auditoria histórica sem perder dados. Timestamps (\texttt{created\_at}, \texttt{updated\_at}) são gerenciados automaticamente
via \texttt{default\_factory}, garantindo consistência sem requerer lógica explícita nos controladores.

Tipos de dados complexos (endereços, itens de pedidos com variações) são serializados como JSON direto no banco,
aproveitando a flexibilidade de plataformas modernas (PostgreSQL, SQLite com extensões) enquanto mantém
queries mais simples e evita JOIN excessivos durante leitura.

Chaves primárias utilizam UUIDs v7, que combinam propriedades de UUIDs tradicionais (unicidade global)
com sortabilidade temporal, favorecendo índices e paginação eficiente.

Os domínios atuais são organizados da seguinte forma:

\begin{verbatim}
├── Audit
│   ├── audit_decorator.py
│   ├── audit_middleware.py
│   ├── audit_repository.py
│   ├── audit_service.py
│   └── audit.py
├── Auth
│   ├── auth_routers.py
│   ├── auth_service.py
│   └── auth.py
├── Catalogue
│   ├── catalogue_repository.py
│   ├── catalogue_routers.py
│   ├── catalogue_service.py
│   ├── catalogue.py
│   └── dish_ingredient.py
├── Costumer
│   ├── costumer_repository.py
│   ├── costumer_routers.py
│   ├── costumer_service.py
│   └── costumer.py
├── Employee
│   ├── employee_repository.py
│   ├── employee_routers.py
│   ├── employee_service.py
│   └── employee.py
├── Inventory
│   ├── inventory_repository.py
│   ├── inventory_routers.py
│   ├── inventory_service.py
│   └── inventory.py
├── Order
│   ├── order_repository.py
│   ├── order_routers.py
│   ├── order_service.py
│   └── order.py
\end{verbatim}
